% Part: computability
% Chapter: computability-theory
% Section: coding-computations

\documentclass[../../../include/open-logic-section]{subfiles}

\begin{document}

\olfileid{cmp}{thy}{cod}
\olsection{د محاسبو کوډول}

د محاسبې په هر مدل کښې لاندې کارونه کېدے شي:
\begin{enumerate}
\item د محاسبه کېدونکو تابعو \emph{تعريفونه} په منظم ډول بيانول۔ د
  بېلګې په توګه، د ټيورينګ ماشين مشخصې، بازګشتي تعريفونه يا د پروګرام
  ليکلو په يوه ژبه کښې پروګرامونه د دې تعريفونو ورکوونکي ګڼلے شئ۔
\item د يو ورکړل شوي ننوت دپاره، د کوم تعريف له مخې د يوې تابع د
  محاسبې بشپړ ريکارډ بيانول۔ د بېلګې په توګه، د ټيورينګ ماشين محاسبه
  د محاسبې په هر ګام کښې د بڼو د لړۍ له لارې بيانېدے شي (د ماشين
  حالت او د ټېپ منځپانګه)۔
\item دا کتل چې د يوې محاسبې يو اټکلي ريکارډ په رښتيا د دې ريکارډ دے
  که نه چې د ورکړل شوي تعريف محاسبه کېدونکې تابع به پر ورکړل شوي ننوت
  څنګه محاسبه شي (هغه ننوت چې تابع پرې تعريف شوې وي، يعنې محاسبه
  درېږي)۔
\item د يوې محاسبې د بشپړ ريکارډ له داسې بيان څخه د ورکړل شوي ننوت
  دپاره د تابع قيمت راايستل۔ د بېلګې په توګه، د يوې درېدونکې ټيورينګ
  ماشين محاسبې په وروستي ګام کښې د ټېپ منځپانګه قيمت دے۔
\end{enumerate}

د کوډولو په کارولو سره د محاسبه کېدونکې تابع هر بيان ته داسې عددي
\emph{شاخص} ټاکل کېدے شي چې لارښوونې په محاسبه کېدونکي ډول له شاخصه
بېرته تر لاسه شي۔ په ورته ډول د يوې محاسبې بشپړ ريکارډ هم په يوه عدد
کوډېدے شي۔ په پايله کښې پيدا شوې حسابي اړيکه، ``$s$ د شاخص~$e$ د تابع
د محاسبې ريکارډ پر ننوت~$x$ کوډوي''، او تابع، ``د کوډ~$s$ لرونکې د
محاسبې لړۍ وت''، بيا محاسبه کېدونکې دي؛ په حقيقت کښې دواړه بنسټيز
بازګشتي دي۔

دا بنسټيز حقيقت ډېر ځواکمن دے، او پرېږدي چې د محاسبه کېدنې په اړه ډېرې
حيرانوونکې او مهمې پايلې د محاسبې له غوره شوي مدل څخه په خپلواک ډول
ثابتې کړو۔

\end{document}
